\section{Каузальность и наблюдаемая скорость света}
\label{sec:c-light-cone}

По решётке сигнал проходит за такт не дальше~$\caSignalSpeed\htick=\caStepSpace$.
На микроуровне мировая линия --- ломаная по рёбрам $N(x)$ со скоростью звена~$\caSignalSpeed$;
массивная частица --- зигзаг внутри того же конуса (\cref{fig:carrier-worldlines-fcc,fig:carrier-worldlines-hex}).
Макроскопический наблюдатель, осредняя поле на масштабах $\gg\caStepSpace$,
восстанавливает null-линию со скоростью $\speedC=\caKappaGeom \caSignalSpeed$ на \latGCK.
Локальная структура размера порядка~$\caStepSpace$ несёт конус~$\caSignalSpeed$;
составная структура, простирающаяся на многие ячейки, усредняется до макроскопической~$c$.
По \cref{prop:m-one-tick-cone,prop:t-speed-limit} на~$\Tlayer$ не возникает фронта
с наблюдаемой скоростью $v>\speedC$; на~$\Mlayer$ сигнал не выходит за конус~$\caSignalSpeed$
(одно ребро за $\htick$).

\begin{figure}[htbp]
  \centering
  \begin{tikzpicture}[carrier panel, x=2.05cm, y=2.05cm]
  \def\tlim{1.2}
  \def\xlim{1.25}
  \def\kappa{0.7071}
  \draw[carrier muted, dashed] (0,0) -- (\xlim,\tlim);
  \draw[carrier muted, dashed] (0,0) -- (-\xlim,\tlim);
  \draw[carrier muted, line width=0.65pt]
    (0,0) -- (0.22,0.18) -- (0.38,0.08) -- (0.55,0.22) -- (0.72,0.12);
  \pgfmathsetmacro{\cmacro}{1.08/\kappa}
  \draw[carrier object] (0,0) -- (\xlim,\xlim/\cmacro);
  \node[anchor=west, font=\scriptsize, inner sep=1pt] at (0.45,0.35) {$\caSignalSpeed$};
  \node[anchor=west, font=\scriptsize, inner sep=1pt] at (0.45,0.72) {$c=\symKappa \caSignalSpeed$};
\end{tikzpicture}

  \caption{Мировые линии на \latGCKshort\ ($\caKappaGeom=1/\sqrt{2}$): микро-null, зигзаг и макро-фотон $\speedC=\caKappaGeom \caSignalSpeed$.}
  \label{fig:carrier-worldlines-fcc}
\end{figure}

\begin{figure}[htbp]
  \centering
  \begin{tikzpicture}[carrier panel, x=2.05cm, y=2.05cm]
  \def\tlim{1.2}
  \def\xlim{1.25}
  \def\kappa{0.8660}
  \draw[carrier muted, dashed] (0,0) -- (\xlim,\tlim);
  \draw[carrier muted, dashed] (0,0) -- (-\xlim,\tlim);
  \draw[carrier muted, line width=0.65pt]
    (0,0) -- (0.22,0.18) -- (0.38,0.08) -- (0.55,0.22) -- (0.72,0.12);
  \pgfmathsetmacro{\cmacro}{1.08/\kappa}
  \draw[carrier object] (0,0) -- (\xlim,\xlim/\cmacro);
  \node[anchor=west, font=\scriptsize, inner sep=1pt] at (0.45,0.35) {$\caSignalSpeed$};
  \node[anchor=west, font=\scriptsize, inner sep=1pt] at (0.45,0.72) {$c=\symKappa \caSignalSpeed$};
\end{tikzpicture}

  \caption{Мировые линии на гекс-срезе $(2{+}1)$ ($\caKappaGeom_{\mathrm{hex}}=\sqrt{3}/2$).}
  \label{fig:carrier-worldlines-hex}
\end{figure}

